\documentclass[a4paper,12pt]{article}
\usepackage[utf8]{inputenc}
\usepackage[german]{babel}
\usepackage{amsmath}
\usepackage[pdftex]{graphicx}
\usepackage{geometry}
\geometry{a4paper, top=15mm, left=15mm, right=15mm, bottom=15mm,
	headsep=10mm, footskip=12mm}
\usepackage{parskip}
\usepackage{href-ul}
\hypersetup{
	colorlinks=true,
	linkcolor=blue,
	urlcolor=blue}

\begin{document}
	
	Worksheet: Simple Electrical Circuits
	
	\section*{Basic Knowledge – Check or Fill in}
	
	\begin{itemize}
		\item \textbf{A Closed Circuit} is necessary for a lamp to \underline{\hspace{3cm}}.
		
		\item The current flows from the \underline{\hspace{2.5cm}} terminal of the battery to the \underline{\hspace{2.5cm}} terminal.
		
		\item In an electrical circuit you need at least:
		\begin{itemize}
			\item ( ) a socket
			\item ( ) a battery or voltage source
			\item ( ) a consumer (e.g., lamp)
			\item ( ) an electric toothbrush
			\item ( ) cables or wires
		\end{itemize}
		
		\item Draw the symbol for:
		\begin{itemize}
			\item Battery: \underline{\hspace{5cm}}
			\item Lamp: \underline{\hspace{5cm}}
			\item Switch: \underline{\hspace{5cm}}
		\end{itemize}
	\end{itemize}
	
	\vspace{1em}
	\hrule
	\vspace{1em}
	
	\section*{Understanding Electrical Circuits }
	
	Aurika builds a simple electrical circuit with a battery, a cable, a switch, and a light bulb.
	
	She closes the switch – and hey: the light bulb lights up!
	
	\begin{itemize}
		\item a) Why does the circuit have to be "closed"?
		
		\vspace{1em}
		\underline{\hspace{12cm}}
		
		\item b) Aurika has two identical light bulbs and a power source.
		What are the possible circuit configurations? Draw two possible circuit diagrams for them!
		What do you need to consider in each case?
		
		\newpage
		\item c) Check:
		In a series circuit of two lamps...
		
		\begin{itemize}
			\item( ) you can turn both lamps on and off independently.
			\item( ) both lamps go out if one burns out.
			\item( ) the same current flows through all lamps, even if they are different.
			\item( ) it gets brighter if you install more lamps.
			\item( ) both lamps must have the same operating current in amperes.
			\item( ) both lamps must have the same operating voltage in volts.
			\item( ) both lamps must have the same power in watts.
			\item( ) the sum of the operating voltages of the lamps must be equal to the voltage of the power source.
		\end{itemize}
		
		\item d) Check:
		In a {\bf parallel circuit } of two lamps...
		
		\begin{itemize}
			\item( ) both lamps can be turned on and off independently of each other.
			\item( ) both lamps go out if one burns out.
			\item( ) the same current flows through all lamps, even if they are different.
			\item( ) it will be brighter if you install more lamps.
			\item( ) both lamps must have the same operating current in amperes.
			\item( ) both lamps must have the same operating voltage in volts.
			\item( ) both lamps must have the same power in watts.
			\item( ) the sum of the operating voltages of the lamps must be equal to the voltage of the power source.
		\end{itemize}
		
		\item e) Now you want to be able to turn the lights in the circuits on and off.
		
		{\bf Series circuit}: Explain why only one switch is sufficient here and draw it in the circuit.
		What happens when the switch is activated?
		
		{\bf Parallel circuit}: Explain why two switches make sense here. Draw these in the circuit.
		What happens when one of the switches is activated? And what happens when both are activated?
		
		\item f) In a household, are electrical devices such as computers, vacuum cleaners, etc. connected to the 230 V mains in parallel or in series?
		
	\end{itemize}
	
	\fbox{
		\begin{minipage}{0.5\textwidth}
			For the solution, please click here or scan the QR code.\\
			More worksheets are available at
			
			\href{https://www.okuyakl.de}{www.okuyakl.de}
		\end{minipage}
		\hfill
		\begin{minipage}{0.4\textwidth}
			\includegraphics[width=1.5 cm]{../../viecher/zwe03}
			\includegraphics[width=3 cm]{qre129}
			\includegraphics[width=2 cm]{../../viecher/afanticon1}
			
	\end{minipage}}
\end{document}